
\subsubsection*{2.2.1}
\phantomsection
\label{sga3:VIIB:2.2.1}

A bialgebra \(\mathbf H\) over \(\mathbf O_k\) is said to be
\emph{flat} if its underlying module is flat
\textup{(cf.\ \SGARefLink{sga3:VIIB:1.2.1}{1.2.1})}.\footnote{\label{sga3:VIIB:nde:79}The editors
have expanded the original in what follows.} If \(\mathbf H\) is
flat, then, by \SGARefLink{sga3:VIIB:1.3.5}{1.3.5},
\[
  A=\Gamma^*(\mathbf H)
\]
is a topologically flat profinite \(k\)-algebra, and
\(\operatorname{Spf}^*(\mathbf H)\), as a functor from
\[
  (\operatorname{Alf}/k)^{\mathrm{op}}
  =
  \operatorname{Vaf}/k
\]
to the category of sets, is isomorphic to the functor
\[
  \operatorname{Spf}(A):
  C\longmapsto
  \operatorname{Hom}_{\operatorname{Vaf}/k}
    \bigl(\operatorname{Spf}(C),\operatorname{Spf}(A)\bigr).
\]
The group structure on \(\operatorname{Spf}^*(\mathbf H)\) therefore
endows
\[
  \mathscr G(\mathbf H)=\operatorname{Spf}(A)
\]
with the structure of a formal group, described explicitly as
follows.

For every object \(C\) of \(\operatorname{Alf}/k\), the \(C\)-module
underlying \(\mathbf H(C)\) is projective. \SGARefLink{sga3:VIIB:1.2.3.A}{Lemma 1.2.3.A} therefore
gives, by induction on \(n\), natural isomorphisms
\[
\begin{aligned}
  \bigl(\mathbf H(C)^*\bigr)^{\widehat\otimes(n+1)}
  &\simeq
  \operatorname{Hom}_C
    \bigl(
      \mathbf H(C),
      \bigl(\mathbf H(C)^*\bigr)^{\widehat\otimes n}
    \bigr)
  \\
  &\simeq
  \operatorname{Hom}_C
    \bigl(
      \mathbf H(C),
      \bigl(\mathbf H(C)^{\otimes n}\bigr)^*
    \bigr)
  \\
  &\simeq
  \bigl(\mathbf H(C)^{\otimes(n+1)}\bigr)^*.
\end{aligned}
\]
For \(n=1,2\), it follows that the \(C\)-algebra structure on
\(\mathbf H(C)\) endows \(\mathbf H(C)^*\) with a diagonal map
satisfying conditions \SGARefLink{sga3:VIIB:2.1}{2.1}\textup{\SGARefLink{sga3:vii:paragraph:5-1:equation:i}{(i)}}--\textup{\SGARefLink{sga3:VIIB:2.1:item:iii}{(iii)}}, functorially
in \(C\).

Consequently,
\[
  A
  =
  \Gamma^*(\mathbf H)
  =
  \varprojlim_{\mathfrak l}\mathbf H(k/\mathfrak l)^*
\]
is endowed with the structure of a cogroup in
\(\operatorname{Alp}/k\), which defines the announced formal-group
structure on \(\mathscr G(\mathbf H)\).

Conversely, let \(G\) be a topologically flat formal \(k\)-group with
affine algebra \(A\), and write
\[
  \mathbf H(G)=\mathbf V_k^f(A)
\]
for its \(\mathbf O_k\)-coalgebra \textup{(cf.\ \SGARefLink{sga3:VIIB:1.2.3}{1.2.3})}. The
diagonal morphism
\[
  \Delta_A:A\longrightarrow A\widehat\otimes_kA
\]
then induces, for every finite-length \(k\)-algebra \(C\), a
\(C\)-linear map
\[
  \mathbf V_k^f(A)(C)
  \otimes_C
  \mathbf V_k^f(A)(C)
  \longrightarrow
  \mathbf V_k^f(A)(C),
\]
which makes the coalgebra \(\mathbf V_k^f(A)(C)\) into a
\(C\)-bialgebra. We shall call \(\mathbf H(G)\) the
\emph{covariant bialgebra} of the formal group \(G\). Thus, by \SGARefLink{sga3:VIIB:1.3.5.D}{Proposition 1.3.5.D}, we obtain the
following.

\medskip
\noindent\textbf{Proposition.}
\phantomsection
\label{sga3:VIIB:2.2.1.prop}
{\itshape
\begin{enumerate}
\item[\textup{(i)}]
The functors
\[
  G\longmapsto\mathbf H(G),
  \qquad
  \mathbf H\longmapsto\mathscr G(\mathbf H)
\]
define an equivalence between the category of topologically flat
formal \(k\)-groups and the category of flat
\(\mathbf O_k\)-bialgebras.

\item[\textup{(ii)}]
This equivalence ``commutes with base change'': if \(k\to K\) is a
morphism of pseudocompact rings, then
\[
  \mathbf H\bigl(G\widehat\otimes_kK\bigr)
  =
  \mathbf H(G)\otimes_kK
\]
and
\[
  \mathscr G\bigl(\mathbf H\otimes_kK\bigr)
  =
  \mathscr G(\mathbf H)\widehat\otimes_kK.
\]
\end{enumerate}}

When \(k\) is an Artinian ring and \(G\) is a topologically flat
formal \(k\)-group, the functor \(\mathbf H(G)\) is plainly
determined by its value
\[
  \mathbf H(G)=\mathbf H(G)(k)
\]
at \(k\). We shall also call \(\mathbf H(G)\) the \emph{(covariant)
bialgebra} of \(G\). Let \(\mathcal H_k\) denote the category of
\(k\)-Hopf algebras, and let
\[
  \mathcal H^{\mathrm{cocom}}_{\mathrm{flat}/k}
\]
be its full subcategory formed by the \(k\)-Hopf algebras that are
flat over \(k\) and cocommutative. We then obtain the following.

\medskip
\noindent\textbf{Corollary.}
\phantomsection
\label{sga3:VIIB:2.2.1.cor}
{\itshape Let \(k\) be an Artinian ring.
\begin{enumerate}
\item[\textup{(i)}]
The functors
\[
  G\longmapsto\mathbf H(G),
  \qquad
  \mathbf H\longmapsto
  \mathscr G(\mathbf H)=\operatorname{Spf}^*(\mathbf H)
\]
define an equivalence between the category of topologically flat
formal \(k\)-groups and
\(\mathcal H^{\mathrm{cocom}}_{\mathrm{flat}/k}\).

\item[\textup{(ii)}]
This equivalence ``commutes with base change'': if \(k\to K\) is a
morphism of Artinian rings, then
\[
  \mathbf H\bigl(G\widehat\otimes_kK\bigr)
  =
  \mathbf H(G)\otimes_kK
\]
and
\[
  \mathscr G\bigl(\mathbf H\otimes_kK\bigr)
  =
  \mathscr G(\mathbf H)\widehat\otimes_kK.
\]
\end{enumerate}}

On the other hand, let
\[
  \mathcal H^{\mathrm{com}}_{\mathrm{flat}/k}
\]
denote the full subcategory of \(\mathcal H_k\) formed by the
\(k\)-Hopf algebras that are flat over \(k\) and commutative. Recall
that the functor
\[
  T\longmapsto\mathcal O(T)
\]
is an anti-equivalence from the category of affine \(k\)-group
schemes onto \(\mathcal H^{\mathrm{com}}_{\mathrm{flat}/k}\).

\subsubsection*{2.2.2}
\phantomsection
\label{sga3:VIIB:2.2.2}

Suppose, for simplicity, that \(k\) is Artinian, and let \(G\) be a
topologically flat formal \(k\)-group.\footnote{%
\label{sga3:VIIB:nde:83}The editors have expanded the original in
what follows in order to introduce the notation \(D'(G)\) and
\(D(K)\); cf.\ [Ca62], \S~14.} Then \(G\) is commutative if and only
if its affine algebra \(\mathcal A(G)\) has a cocommutative
comultiplication, which is equivalent to saying that the bialgebra
\(\mathbf H(G)\) has a commutative multiplication. In that case,
\(\mathbf H(G)\) is a commutative and cocommutative Hopf algebra,
flat over \(k\). Thus, if
\[
  D'(G)=\operatorname{Spec}\mathbf H(G),
\]
then \(D'(G)\) is a commutative \(k\)-group scheme, affine and flat
over \(k\).

Conversely, if \(T\) is such a \(k\)-group scheme, its affine algebra
\(\mathcal O(T)\) is a commutative group object in the category of
cocommutative \(k\)-coalgebras flat over \(k\). Hence \SGARefLink{sga3:VIIB:1.3.5.C}{1.3.5.C} gives
a topologically flat formal \(k\)-group \(D(T)\), defined by
\[
  D(T)
  =
  \operatorname{Spf}^*\mathcal O(T)
  =
  \operatorname{Spf}(\mathcal A),
  \qquad
  \mathcal A=\mathcal O(T)^*.
\]
By \SGARefLink{sga3:VIIB:1.3.5.D}{1.3.5.D}, there are canonical isomorphisms
\[
  G=\operatorname{Spf}^*\mathbf H(G),
  \qquad
  \mathbf H(D(T))=\mathcal O(T).
\]
We therefore obtain canonical isomorphisms
\[
\begin{aligned}
  D(D'(G))
  &=
  \operatorname{Spf}^*\mathcal O(D'(G))
  =
  \operatorname{Spf}^*\mathbf H(G)
  =
  G,
  \\
  D'(D(T))
  &=
  \operatorname{Spec}\mathbf H(D(T))
  =
  \operatorname{Spec}\mathcal O(T)
  =
  T.
\end{aligned}
\]
Moreover, if \(k\text{-}\operatorname{Gr.}\) denotes the category of
\(k\)-group schemes, \SGARefLink{sga3:VIIB:2.2.1}{Corollary 2.2.1} gives functorial isomorphisms
\[
  \operatorname{Hom}_{\operatorname{Grf}/k}\bigl(G,D(T)\bigr)
  \simeq
  \operatorname{Hom}_{\mathcal H_k}
    \bigl(\mathbf H(G),\mathcal O(T)\bigr)
  \simeq
  \operatorname{Hom}_{k\text{-}\operatorname{Gr.}}
    \bigl(T,D'(G)\bigr).
\]
This proves the following.

\medskip
\noindent\textbf{Proposition (Cartier duality).}
\phantomsection
\label{sga3:VIIB:2.2.2.prop}
{\itshape Let \(k\) be an Artinian ring.
\begin{enumerate}
\item[\textup{(i)}]
The functors
\[
  G\longmapsto D'(G)=\operatorname{Spec}\mathbf H(G),
  \qquad
  T\longmapsto D(T)=\operatorname{Spf}^*\mathcal O(T)
\]
induce an anti-equivalence between the category
\(\mathcal{FC}_k\) of commutative topologically flat formal
\(k\)-groups and the category \(\mathcal{AC}_k\) of commutative
affine flat \(k\)-group schemes. In other words, \(G\) and \(D'(G)\),
respectively \(T\) and \(D(T)\), are related by the equalities
\footnote{\label{sga3:VIIB:nde:84}\textit{Editors' note.} Let
\(\mathcal{FC}_k^f\), respectively \(\mathcal{AC}_k^f\), denote the
full subcategory of \(\mathcal{FC}_k\), respectively
\(\mathcal{AC}_k\), consisting of the objects \(G\), respectively
\(T\), for which \(\mathbf H(G)\), respectively \(\mathcal O(T)\),
is a finite \(k\)-module \textup{(and hence finite locally free)}.
Then \(\mathcal{FC}_k^f\) and \(\mathcal{AC}_k^f\) both have the
same objects as the category \(\mathcal C\) of commutative and
cocommutative Hopf \(k\)-algebras that are finite and flat over
\(k\). The correspondence \(G\mapsto\mathbf H(G)\), respectively
\(T\mapsto\mathcal O(T)\), is covariant, respectively
contravariant. Thus one recovers the ``Cartier duality'' of the
category \(\mathcal{AC}_k\), already seen in
VII\textsubscript{A} \SGARefLink{sga3:VIIA:3.3.1}{3.3.1}.}
\[
  \mathbf H(G)=\mathcal O(D'(G)),
  \qquad
  \mathcal O(T)=\mathbf H(D(T)).
\]

\item[\textup{(ii)}]
This anti-equivalence ``commutes with base change'': if \(k\to K\)
is a morphism of Artinian rings, then
\[
  D'\bigl(G\widehat\otimes_kK\bigr)
  =
  D'(G)\otimes_kK
\]
and
\[
  D\bigl(T\otimes_kK\bigr)
  =
  D(T)\widehat\otimes_kK.
\]

\item[\textup{(iii)}]
In particular, if \(k\) is a field, one obtains an
anti-equivalence between the category of commutative formal
\(k\)-groups and the category of commutative affine \(k\)-group
schemes, and this anti-equivalence commutes with extension of the
base field.
\end{enumerate}}

